\documentclass[11pt,a4paper]{article}
\usepackage[utf8]{inputenc}
\usepackage[brazil]{babel}
\usepackage{amsmath,amssymb}
\usepackage{booktabs}
\usepackage{geometry}
\usepackage{enumitem}
\usepackage{array}
\usepackage{longtable}

\geometry{margin=2.5cm}

\title{Formulação Matemática do Modelo de Otimização\\de Alocação de Produção de Ovos}
\author{Mantiqueira}
\date{\today}

\begin{document}

\maketitle

\tableofcontents
\newpage

%==============================================================================
\section{Visão Geral}
%==============================================================================

O modelo de otimização tem como objetivo maximizar a margem total de contribuição na alocação da produção semanal de ovos para diferentes SKUs (Stock Keeping Units), respeitando:

\begin{itemize}
    \item Capacidade de produção por classe de produto
    \item Demanda histórica como limite superior por SKU
    \item Reservas prioritárias para pedidos especiais (GRANEL, Exportação, Marca Própria)
    \item Restrições de integralidade (caixas completas)
\end{itemize}

%==============================================================================
\section{Conjuntos e Índices}
%==============================================================================

\subsection{Conjuntos Principais}

\begin{align}
I &= \{1, 2, \ldots, n\} && \text{Conjunto de todos os SKUs} \\
C &= \{c_1, c_2, \ldots, c_m\} && \text{Conjunto de classes de produtos} \\
S &= \{s_1, s_2, \ldots, s_k\} && \text{Conjunto de semanas no histórico} \\
F &= \{f_1, f_2, \ldots, f_l\} && \text{Conjunto de registros de faturamento}
\end{align}

\subsection{Subconjuntos Derivados}

\begin{align}
I_c &\subseteq I && \text{SKUs pertencentes à classe } c \\
I_{\text{válido}} &\subseteq I && \text{SKUs válidos para otimização} \\
I_{\text{prod}} &\subseteq I && \text{SKUs com produção na semana corrente} \\
H &\subseteq I_{\text{válido}} && \text{SKUs com demanda histórica válida} \\
R &\subseteq I && \text{SKUs com pedidos reservados}
\end{align}

\subsection{Mapeamentos}

\begin{align}
\kappa: I &\rightarrow C && \text{Função que mapeia SKU para sua classe} \\
\varepsilon: I &\rightarrow \mathbb{Z}^+ && \text{Função que retorna ovos por embalagem} \\
\sigma: I &\rightarrow S && \text{Função que mapeia SKU para semanas com venda}
\end{align}

%==============================================================================
\section{Parâmetros de Entrada}
%==============================================================================

\subsection{Dados do SKU}

\begin{table}[h]
\centering
\begin{tabular}{clll}
\toprule
\textbf{Símbolo} & \textbf{Descrição} & \textbf{Unidade} & \textbf{Domínio} \\
\midrule
$p_i$ & Preço de venda do SKU $i$ & R\$/caixa & $\mathbb{R}^+$ \\
$c_i$ & Custo YTD do SKU $i$ & R\$/caixa & $\mathbb{R}^+$ \\
$e_i$ & Quantidade de ovos por caixa & ovos/caixa & $\mathbb{Z}^+$ \\
$\kappa(i)$ & Classe do SKU $i$ & - & $C$ \\
$d_i$ & Descrição do SKU $i$ & - & texto \\
\bottomrule
\end{tabular}
\caption{Parâmetros por SKU}
\end{table}

\subsection{Cálculo da Embalagem}

A quantidade de ovos por embalagem $e_i$ é extraída da descrição do produto:

\begin{equation}
e_i = n_{\text{bandejas}} \times n_{\text{unidades}}
\end{equation}

Exemplo: ``CX 12 BJ 30 UN'' $\Rightarrow e_i = 12 \times 30 = 360$ ovos

\subsection{Dados de Produção}

\begin{align}
P_c &\in \mathbb{R}^+ && \text{Produção total da classe } c \text{ na semana (ovos)} \\
P_c &= \sum_{i \in I_c \cap I_{\text{prod}}} \hat{p}_i \cdot e_i && \text{onde } \hat{p}_i \text{ é produção em caixas}
\end{align}

\subsection{Margem de Contribuição}

\textbf{Margem por caixa:}
\begin{equation}
m_i = p_i - c_i \quad \text{(R\$/caixa)}
\end{equation}

\textbf{Margem por ovo:}
\begin{equation}
\mu_i = \frac{m_i}{e_i} = \frac{p_i - c_i}{e_i} \quad \text{(R\$/ovo)}
\end{equation}

\textbf{Tratamento de custo ausente:}

Quando um SKU $i$ não possui custo individual, utiliza-se a média da classe:
\begin{equation}
c_i = \begin{cases}
c_i^{\text{original}} & \text{se } c_i^{\text{original}} > 0 \\[4pt]
\displaystyle\frac{\sum_{j \in I_{\kappa(i)}} c_j}{|I_{\kappa(i)}|} & \text{caso contrário}
\end{cases}
\end{equation}

%==============================================================================
\section{Processamento de Dados (ETL)}
%==============================================================================

\subsection{Base de Faturamento}

Cada registro de faturamento $f \in F$ contém:
\begin{equation}
f = (\text{item}_f, \text{cliente}_f, \text{estab}_f, \text{data}_f, q_f, \text{desc}_f)
\end{equation}

onde $q_f$ representa a quantidade em \textbf{caixas equivalentes de 360 ovos} (base normalizada).

\subsection{Correção de Estabelecimento}

Alguns registros de faturamento têm estabelecimento incorreto devido a transferências internas. A correção é feita via mapeamento:

\begin{equation}
E: \text{Cod.Cliente} \rightarrow \text{Estab\_Padrao}
\end{equation}

\begin{equation}
\text{Estab\_Corrigido}(f) = \begin{cases}
E(\text{cliente}_f) & \text{se } \text{cliente}_f \in \text{dom}(E) \\
\text{estab}_f & \text{caso contrário}
\end{cases}
\end{equation}

\subsection{Filtro por Estabelecimento}

Apenas registros do estabelecimento de interesse são considerados:
\begin{equation}
F_{\text{filtrado}} = \{f \in F : \text{Estab\_Corrigido}(f) = 100\}
\end{equation}

\subsection{Agregação Temporal}

Para cada registro $f$, extrai-se a semana ISO:
\begin{equation}
\text{semana}(f) = \text{ISO\_Week}(\text{data}_f) = (y, w) \quad \text{onde } y = \text{ano}, w = \text{semana}
\end{equation}

\subsection{Cálculo da Demanda Histórica}

\subsubsection{Quantidade Semanal por SKU}

Para cada SKU $i$ e semana $s$:
\begin{equation}
Q_{i,s} = \sum_{f \in F : \text{item}_f = i, \text{semana}(f) = s} q_f \times 360 \quad \text{(ovos)}
\end{equation}

O fator 360 converte caixas normalizadas para quantidade de ovos.

\subsubsection{Limite de Demanda Histórica}

\textbf{Método Máximo:}
\begin{equation}
D_i^{\max} = \max_{s \in S} Q_{i,s}
\end{equation}

\textbf{Método Média:}
\begin{equation}
D_i^{\text{média}} = \frac{1}{|S_i|} \sum_{s \in S_i} Q_{i,s}
\end{equation}

onde $S_i = \{s \in S : Q_{i,s} > 0\}$ é o conjunto de semanas com vendas do SKU $i$.

\textbf{Parâmetro de configuração:}
\begin{equation}
D_i = \begin{cases}
D_i^{\max} & \text{se } \texttt{tipo\_calculo\_demanda} = \texttt{maximo} \\
D_i^{\text{média}} & \text{se } \texttt{tipo\_calculo\_demanda} = \texttt{media}
\end{cases}
\end{equation}

\subsection{Classificação de SKUs Válidos}

Um SKU é \textbf{válido para otimização} se atende todos os critérios:

\begin{equation}
i \in I_{\text{válido}} \iff \begin{cases}
\text{ESTAB}_i = 100 & \text{(estabelecimento correto)} \\
\text{STATUS}_i = \texttt{ATIVO} & \text{(produto ativo)} \\
\text{TIPO}_i \notin T_{\text{reserva}} & \text{(tipo permitido)} \\
e_i > 0 & \text{(embalagem válida)}
\end{cases}
\end{equation}

onde $T_{\text{reserva}} = \{\texttt{GRANEL}, \texttt{Exportação}, \texttt{MARCA PROPRIA}\}$

\subsection{Conjunto com Demanda Histórica}

\begin{equation}
H = \{i \in I_{\text{válido}} : D_i > 0\}
\end{equation}

SKUs em $I_{\text{válido}} \setminus H$ têm produção mas não participam da otimização por falta de validação histórica.

%==============================================================================
\section{Sistema de Reservas de Pedidos}
%==============================================================================

\subsection{Origem das Reservas}

Reservas são geradas para SKUs de tipos especiais com base na demanda histórica média:

\begin{equation}
Q_i = \frac{1}{|S_i|} \sum_{s \in S_i} Q_{i,s} \quad \forall i : \text{TIPO}_i \in T_{\text{reserva}}
\end{equation}

\subsection{Critério de Elegibilidade para Reserva}

Um pedido para SKU $i$ é \textbf{elegível para reserva} se:

\begin{equation}
i \in R \iff \begin{cases}
Q_i > 0 & \text{(existe pedido)} \\
i \in I_{\text{prod}} & \text{(tem produção na semana)} \\
e_i > 0 & \text{(embalagem válida/extraível)}
\end{cases}
\end{equation}

\subsection{Cálculo da Quantidade Reservada}

Para cada SKU $i \in R$:
\begin{equation}
r_i = \min(Q_i, P_{\kappa(i)})
\end{equation}

A reserva é limitada pela produção disponível da classe.

\subsection{Agregação de Reservas por Classe}

Total reservado por classe:
\begin{equation}
R_c = \sum_{i \in R : \kappa(i) = c} r_i \quad \forall c \in C
\end{equation}

\subsection{Produção Excedente (Disponível para Otimização)}

\begin{equation}
E_c = \max(0, P_c - R_c) \quad \forall c \in C
\end{equation}

\subsection{Pedidos Ignorados}

Um pedido é \textbf{ignorado} (não reservado) se:

\begin{equation}
i \in I_{\text{ignorado}} \iff (Q_i > 0) \land \left[(i \notin I_{\text{prod}}) \lor (e_i = 0)\right]
\end{equation}

\textbf{Motivos de ignorar:}
\begin{itemize}
    \item SKU sem produção na semana corrente
    \item Embalagem não extraível da descrição (ex: kits)
    \item Classe não mapeada
\end{itemize}

%==============================================================================
\section{Modelo de Otimização}
%==============================================================================

\subsection{Variável de Decisão}

\textbf{Implementação atual (OR-Tools):} a variável é \textbf{contínua} e expressa em ovos:

\begin{equation}
x_i \in \mathbb{R}^+ \quad \text{Quantidade alocada ao SKU } i \text{ (em ovos)}
\end{equation}

O solver utiliza \texttt{NumVar} (variável contínua). O resultado é tratado/arredondado posteriormente para fins de relatório.

\textbf{Formulação teórica (caixas inteiras):} se fosse exigido que a alocação corresponda apenas a caixas inteiras, teríamos $x_i \in \mathbb{Z}^+$ com $x_i \equiv 0 \pmod{e_i}$, ou equivalentemente uma variável $y_i \in \mathbb{Z}^+$ em caixas (ver Seção ``Reformulação em Caixas''). Essa formulação inteira \textbf{não está implementada} na versão atual.

\subsection{Função Objetivo}

Maximizar a margem total de contribuição:

\begin{equation}
\max Z = \sum_{i \in H} \mu_i \cdot x_i = \sum_{i \in H} \frac{p_i - c_i}{e_i} \cdot x_i
\end{equation}

\textbf{Interpretação:} Selecionar SKUs com maior margem por ovo, respeitando as restrições.

\subsection{Restrições do Modelo}

\subsubsection{R1: Capacidade de Produção por Classe}

A soma das alocações de uma classe não pode exceder a produção excedente:

\begin{equation}
\sum_{i \in H : \kappa(i) = c} x_i \leq E_c \quad \forall c \in C
\end{equation}

\textbf{Interpretação:} Após descontar as reservas, o modelo só pode alocar o que sobrou.

\subsubsection{R2: Limite de Demanda Histórica}

A alocação individual é limitada pela demanda histórica:

\begin{equation}
x_i \leq D_i \quad \forall i \in H
\end{equation}

\textbf{Interpretação:} Evita alocações para SKUs além do que o mercado historicamente absorve.

\subsubsection{R3: Não-Negatividade}

\begin{equation}
x_i \geq 0 \quad \forall i \in H
\end{equation}

\subsubsection{R4: Integralidade por Embalagem (não implementada)}

Na formulação teórica, a alocação deveria corresponder a caixas inteiras:

\begin{equation}
x_i = k_i \cdot e_i \quad \text{onde } k_i \in \mathbb{Z}^+ \quad \forall i \in H
\end{equation}

Equivalentemente: $x_i \equiv 0 \pmod{e_i}$. Na \textbf{implementação atual} essa restrição não é aplicada: as variáveis são contínuas ($x_i \in \mathbb{R}^+$).

%==============================================================================
\section{Formulação Compacta}
%==============================================================================

\textbf{Implementação atual (variáveis contínuas em ovos):}

\begin{equation}
\boxed{
\begin{aligned}
\max \quad & Z = \sum_{i \in H} \mu_i \cdot x_i \\[10pt]
\text{sujeito a:} \quad & \\[4pt]
& \sum_{i \in H : \kappa(i) = c} x_i \leq E_c && \forall c \in C \quad \text{(capacidade)} \\[6pt]
& x_i \leq D_i && \forall i \in H \quad \text{(demanda histórica)} \\[6pt]
& x_i \geq 0 && \forall i \in H \quad \text{(não-negatividade)} \\[6pt]
& x_i \in \mathbb{R}^+ && \forall i \in H \quad \text{(contínua)}
\end{aligned}
}
\end{equation}

\textbf{Formulação teórica com caixas inteiras} (não implementada): acrescentar $x_i = k_i \cdot e_i$, $k_i \in \mathbb{Z}^+$, ou usar a reformulação em $y_i$ (Seção seguinte).

%==============================================================================
\section{Reformulação em Caixas (Alternativa)}
%==============================================================================

\textbf{Nota:} A implementação atual usa variáveis contínuas em \textbf{ovos} (Seção anterior). A reformulação abaixo é uma \textbf{alternativa} para o caso em que se deseja variáveis \textbf{inteiras em caixas}.

Se a variável for definida em termos de caixas:

\begin{equation}
y_i = \frac{x_i}{e_i} \in \mathbb{Z}^+ \quad \text{(quantidade em caixas)}
\end{equation}

obtém-se o modelo reformulado (equivalente, com integralidade):

\begin{equation}
\begin{aligned}
\max \quad & Z = \sum_{i \in H} m_i \cdot y_i \\[8pt]
\text{s.a.} \quad & \sum_{i \in H : \kappa(i) = c} e_i \cdot y_i \leq E_c && \forall c \in C \\[4pt]
& e_i \cdot y_i \leq D_i && \forall i \in H \\[4pt]
& y_i \in \mathbb{Z}^+ && \forall i \in H
\end{aligned}
\end{equation}

Esta formulação em $y_i$ \textbf{não está implementada} no código atual; o solver utiliza $x_i \in \mathbb{R}^+$ (ovos, contínuo).

%==============================================================================
\section{Fluxo de Execução Detalhado}
%==============================================================================

\begin{enumerate}
    \item \textbf{Carregar configuração} (\texttt{config.yaml})
    \begin{itemize}
        \item Caminhos dos arquivos de entrada
        \item Parâmetros: semana de referência, tipo de cálculo de demanda
        \item Flags: atender pedidos, usar apenas excedente
    \end{itemize}
    
    \item \textbf{Carregar dados brutos}
    \begin{itemize}
        \item Produção: \texttt{PRODUÇÃO DIA.xlsx}
        \item Custos: \texttt{CUSTO ITEM.csv}
        \item Preços: \texttt{precos\_sku\_embalagem.csv}
        \item Classes: \texttt{base\_skus\_classes.xlsx}
        \item Faturamento: \texttt{manti\_fat\_2025\_full.parquet}
        \item SKUs restritos: \texttt{skus\_restritos.xlsx}
        \item Pedidos: \texttt{pedidos\_clientes.csv}
    \end{itemize}
    
    \item \textbf{Aplicar correção de estabelecimento}
    \begin{itemize}
        \item Carregar mapeamento \texttt{ESTAB CORRIGIDO.xlsx}
        \item Criar coluna \texttt{Estab\_Corrigido}
    \end{itemize}
    
    \item \textbf{Filtrar SKUs válidos}
    \begin{itemize}
        \item ESTAB = 100
        \item STATUS = ATIVO
        \item TIPO $\notin$ \{GRANEL, Exportação, MARCA PROPRIA\}
    \end{itemize}
    
    \item \textbf{Calcular demanda histórica}
    \begin{itemize}
        \item Filtrar faturamento por Estab\_Corrigido = 100
        \item Agregar por SKU e semana
        \item Aplicar máximo ou média conforme configuração
    \end{itemize}
    
    \item \textbf{Processar pedidos/reservas}
    \begin{itemize}
        \item Verificar elegibilidade (produção na semana, embalagem válida)
        \item Calcular reservas $r_i$
        \item Identificar pedidos ignorados
    \end{itemize}
    
    \item \textbf{Calcular produção excedente}
    \begin{itemize}
        \item $E_c = P_c - R_c$ por classe
    \end{itemize}
    
    \item \textbf{Montar e resolver modelo de otimização}
    \begin{itemize}
        \item Criar variáveis de decisão
        \item Adicionar restrições
        \item Executar solver (OR-Tools CP-SAT)
    \end{itemize}
    
    \item \textbf{Extrair e salvar resultados}
    \begin{itemize}
        \item Alocação otimizada por SKU
        \item Reservas por SKU
        \item Métricas financeiras
        \item Pedidos ignorados
    \end{itemize}
\end{enumerate}

%==============================================================================
\section{Métricas de Saída}
%==============================================================================

\subsection{Métricas de Volume}

\begin{align}
\text{Quantidade Alocada Total} &= \sum_{i \in H} x_i^* \\[6pt]
\text{Quantidade Reservada Total} &= \sum_{i \in R} r_i \\[6pt]
\text{Quantidade Destinada Total} &= \sum_{i \in H} x_i^* + \sum_{i \in R} r_i \\[6pt]
\text{Saldo Não Alocado} &= \sum_{c \in C} P_c - \text{Quantidade Destinada Total}
\end{align}

\subsection{Métricas Financeiras}

\begin{align}
\text{Margem Total Otimização} &= \sum_{i \in H} \mu_i \cdot x_i^* = Z^* \\[6pt]
\text{Receita Total} &= \sum_{i \in H} \frac{p_i}{e_i} \cdot x_i^* \\[6pt]
\text{Custo Total} &= \sum_{i \in H} \frac{c_i}{e_i} \cdot x_i^*
\end{align}

\subsection{Métricas de Eficiência}

\begin{align}
\text{Utilização da Produção} &= \frac{\text{Quantidade Destinada Total}}{\sum_{c \in C} P_c} \times 100\% \\[6pt]
\text{Taxa de Atendimento de Pedidos} &= \frac{|R|}{|R| + |I_{\text{ignorado}}|} \times 100\%
\end{align}

%==============================================================================
\section{Glossário de Símbolos}
%==============================================================================

\begin{longtable}{cl}
\toprule
\textbf{Símbolo} & \textbf{Descrição} \\
\midrule
\endhead
$I$ & Conjunto de todos os SKUs \\
$C$ & Conjunto de classes de produtos \\
$H$ & Conjunto de SKUs com demanda histórica válida \\
$R$ & Conjunto de SKUs com reservas \\
$i$ & Índice de SKU \\
$c$ & Índice de classe \\
$s$ & Índice de semana \\
$p_i$ & Preço do SKU $i$ (R\$/caixa) \\
$c_i$ & Custo do SKU $i$ (R\$/caixa) \\
$e_i$ & Ovos por embalagem do SKU $i$ \\
$m_i$ & Margem por caixa do SKU $i$ \\
$\mu_i$ & Margem por ovo do SKU $i$ \\
$P_c$ & Produção total da classe $c$ (ovos) \\
$D_i$ & Limite de demanda histórica do SKU $i$ (ovos) \\
$Q_i$ & Quantidade de pedido/reserva do SKU $i$ (ovos) \\
$r_i$ & Quantidade reservada do SKU $i$ (ovos) \\
$R_c$ & Total reservado da classe $c$ (ovos) \\
$E_c$ & Produção excedente da classe $c$ (ovos) \\
$x_i$ & Variável de decisão: alocação do SKU $i$ (ovos); na implementação atual $x_i \in \mathbb{R}^+$ (contínua) \\
$y_i$ & Variável auxiliar: alocação do SKU $i$ (caixas) \\
$\kappa(i)$ & Classe do SKU $i$ \\
$Z$ & Função objetivo (margem total) \\
\bottomrule
\end{longtable}

\end{document}
